\section{A polarized degeneration with the prescribed isotropic class}
\label{sec:family}

The degeneration must supply two kinds of data. Its volume growth and
a nonzero residue at a maximal crossing will allow us to apply
Theorem~\ref{thm:torus}. Its marked periods and second-cohomology
monodromy will retain the prescribed isotropic class for rotation.
We construct both in the same polarized family.
All markings take values in
\[
 \Lambda=H^2(X,\Z),\qquad \Lambda_{\Q}=\Lambda\otimes_{\Z}\Q,
 \qquad r=\operatorname{rank}\Lambda,
\]
with its integral Beauville--Bogomolov--Fujiki form $q$.  We write
$q(x,y)$ for the associated symmetric bilinear form, so that
$q(x)=q(x,x)$.  Markings are obtained by deformation and parallel
transport from the differentiable manifold underlying $X$.
Fix the connected Teichm\"uller component containing $X$. Along
deformation paths use Ehresmann trivializations anchored at $X$;
the transported complex structures on its differentiable manifold
then stay in this component.

The universal coefficient theorem and $H_1(X,\Z)=0$ show that
$\Lambda$ is torsion-free.  Write $c_1(L)=m_0e$ with
$m_0\in\Z_{>0}$ and $e\in\Lambda$ primitive.  Since
$H^1(X,\mathcal O_X)=0$ by Hodge theory, the
exponential sequence identifies $\Pic(X)$ with the integral
$(1,1)$ classes.  Thus $e=c_1(E)$ for a unique line bundle $E$,
and $L\simeq E^{\otimes m_0}$.  The class $e$ is nonzero, nef,
and isotropic.  It will remain the same labelled lattice vector,
although it need not remain of type $(1,1)$ during the construction.
Its sign is fixed by
\[
                    q(e,\kappa_X)>0
\]
for every Kähler class $\kappa_X$ on $X$: in the Lorentzian space
$H^{1,1}(X,\R)$, a nonzero null vector in the closure of the positive
cone pairs positively with every vector in that cone.

\begin{samepage}
\begin{proposition}[The prescribed-cusp family]
\label{prop:family}\label{prop:polarized-family}
Put $d=2n$.  There exist a primitive vector $h\in\Lambda$ with
$q(h)>0$ and $q(h,e)=0$, a rational subspace
\[
 P=\langle e,f,v\rangle_{\Q}\subset h^\perp,
 \qquad \operatorname{sign}(q|_P)=(2,1),
\]
and a projective semistable family $\pi:\mathcal Y\to\Delta$ with
the following properties.  Here $\Delta$ is a sufficiently small disk
in a smooth algebraic curve, and $\Delta^*=\Delta\setminus\{0\}$.
\begin{enumerate}
\item The total space $\mathcal Y$ is smooth, the central fiber is a
reduced simple normal crossings divisor, and every $Y_t$, $t\ne0$,
is an irreducible holomorphic symplectic manifold in the marked
deformation component of $X$.  A line bundle $\mathcal A$ on
$\mathcal Y$ restricts to an ample bundle on every $Y_t$ with
$c_1(\mathcal A|_{Y_t})=kh$ for a fixed positive integer $k$ under
parallel transport.  The family is the analytic restriction of an
algebraic family.
\item The rational vectors can be chosen so that
\begin{equation}
 \label{eq:family-hyperbolic-basis}
 q(e)=q(f)=0,\quad q(e,f)=1,\quad
 q(e,v)=q(f,v)=0,\quad a:=q(v)>0.
\end{equation}
On the universal cover of $\Delta^*$ the marked Hodge line is spanned
by the normalized vector
\begin{equation}
 \label{eq:family-period}
 p(\tau)=f+\tau v-\frac a2\tau^2e,
 \qquad \operatorname{Im}\tau>0.
\end{equation}
The function $\tau$ tends to the cusp represented by the original
line $\Q e$, and
\begin{equation}
 \label{eq:family-period-identities}
 q(p(\tau))=0,\qquad q(p(\tau),e)=1,\qquad
 q(p(\tau),\overline{p(\tau)})
       =2a(\operatorname{Im}\tau)^2.
\end{equation}
For a consistent convention for second-cohomology monodromy there
is a rational $w>0$ such that
\begin{equation}
 \label{eq:family-monodromy}
 Me=e,\qquad Mv=v-aw e,\qquad
 Mf=f+w v-\frac{aw^2}{2}e,
 \qquad M|_{P^\perp}=1.
\end{equation}
The orthogonal complement in the last formula is taken in
$\Lambda_{\Q}$.  In particular,
\begin{equation}
 \label{eq:family-monodromy-image}
 (M-1)P=\langle e,v\rangle_{\Q},\quad
 (M-1)^2f=-aw^2e\ne0,\quad (M-1)^3=0.
\end{equation}
The same image statement holds if the monodromy convention uses
$M^{-1}$.
\item There is a single-valued holomorphic relative symplectic form
$\sigma_t$ on the smooth family, uniquely normalized by
$q([\sigma_t],e)=1$.  The forms $\eta_t=\sigma_t^n$ extend to a
holomorphic section of the relative logarithmic top-form line
bundle of the semistable family.  With $s=-\log|t|$,
\begin{equation}
 \label{eq:family-volume-growth}
 c s^d\le \int_{Y_t}|\eta_t|^2\le C s^d
\end{equation}
for positive constants $c,C$.  There is a point $o\in Y_0$ where
$d+1$ components meet and coordinates satisfying
\begin{equation}
 \label{eq:family-maximal-residue}
 t=z_0\cdots z_d,\qquad
 \eta=F(z)\,\frac{\dd z_1}{z_1}\wedge\cdots\wedge
                   \frac{\dd z_d}{z_d},\qquad F(o)\ne0.
\end{equation}
\item For any prescribed irrational number $\lambda$, there is a
path tending to $0$ in $\Delta^*$, with a continuous choice of lift,
along which $\tau=\lambda+iy$ and $y\to+\infty$.
\end{enumerate}
\end{proposition}
\end{samepage}

The prescribed line $\Q e$ and the normalization $q(p,e)=1$ are the
data needed for rotation.  We obtain them in three stages.  First,
choose a positive polarization perpendicular to $e$ and realize an
open set of its periods by an algebraic family.  Next, find a rational
three-dimensional period subspace containing $e$ and follow its cusp
inside that family, retaining the original marking.  Finally, compute
the periods at the cusp and use their growth to produce the nonzero
logarithmic residue required by Theorem~\ref{thm:torus}.

\subsection{A positive polarization orthogonal to the null class}

We use the following standard inputs from hyperkähler period theory:
unobstructed deformations and local Torelli; surjectivity of the
period map from the selected connected marked component; the
projectivity criterion for a positive integral $(1,1)$ class; and
the Kähler-cone description, whose walls are orthogonal to negative
integral $(1,1)$ classes.  See
\cite{Huybrechts,HuybrechtsErratum,Verbitsky,VerbitskyErratum}.
The corrected projectivity criterion is
\cite[Theorem~2]{HuybrechtsErratum}, and the K\"ahler-cone description
in terms of negative MBM classes is due to Amerik--Verbitsky
\cite[Theorems~1.19 and~6.2]{AmerikVerbitsky}; see also
\cite[Section~2.1]{SV}.

The first two construction steps also have an independent use: they
provide an algebraic family and a quotient in which equality of period
lifts detects equality of monodromy on the full second-cohomology
lattice. We isolate precisely this construction, before choosing a
ternary subdomain or a cusp.

\begin{lemma}[A positive polarization and a free period quotient]
\label{lem:polarized-family-construction}
Let $X$ be a compact irreducible holomorphic symplectic K\"ahler
manifold, fix its connected marked deformation component, and put
$\Lambda=H^2(X,\Z)$ with its integral BBF form $q$.
For each nonzero primitive $e\in\Lambda$ with $q(e)=0$, there are
\begin{enumerate}
\item a primitive $h\in\Lambda$ with $q(h)>0$ and $q(h,e)=0$,
and a manifold in the chosen marked component on which $h$ is ample;
\item a smooth projective algebraic family over a smooth connected
quasi-projective base $S$, whose fibers are irreducible holomorphic
symplectic manifolds in that component, with a relative polarization
of marked class $kh$ for one $k\in\Z_{>0}$;
\item a torsion-free finite-index subgroup
\[
 \Gamma\subset G_h(\Z)
 :=\{g\in O(\Lambda,q):gh=h,\ gD_h=D_h\},
\]
where $D_h$ is the component of the polarized period domain containing
the chosen ample period, such that the family monodromy is contained
in $\Gamma$ and the descended period map
\[
       \Phi:S\longrightarrow Q_h:=\Gamma\backslash D_h
\]
is algebraic. The quotient $Q_h$ is smooth and quasi-projective,
$D_h\to Q_h$ is a covering map with free $\Gamma$-action, and a local
marked lift of $\Phi$ is submersive at some point of $S$.
\end{enumerate}
More generally, conclusions \textup{(2)} and \textup{(3)} hold for any
prescribed primitive positive $h\in\Lambda$ which is ample on a
manifold in the chosen marked component.
The markings are obtained by deformation and parallel transport from
$X$. No finite-index assertion about the image of the family monodromy
inside $\Gamma$ is part of this construction.
\end{lemma}

\begin{proof}
\emph{The positive polarization.}
The signature of $q$ is $(3,r-3)$.  The existence of the nonzero
null vector $e$ implies $r\ge4$.  The radical of $e^\perp$ is $\R e$,
and
\[
 \operatorname{sign}(e^\perp/\R e)=(2,r-4).
\]
Thus $e^\perp$ contains positive real vectors.  Since $e$ is rational,
$e^\perp$ is defined over $\Q$; rational approximation inside this
space supplies a rational positive vector. Taking the primitive
integral generator of its ray gives $h$ with
$q(h)>0$ and $q(h,e)=0$.

Set $V=h^\perp\subset\Lambda_{\Q}$; its signature is $(2,r-3)$.
The polarized period domain is
\[
 \{[p]\in\mathbb P(V_{\C}):q(p)=0,\ q(p,\bar p)>0\}.
\]
A generic point of this domain is orthogonal to no integral class
outside $\Q h$: the excluded sets are countably many proper
hyperplane sections.  Period surjectivity realizes such a point
in the selected marked component.  Its Néron--Severi lattice is
$\Z h$.  One of $h,-h$ is in the component of the positive cone
containing the Kähler cone.  There are no negative integral
$(1,1)$ classes on this manifold, so the Kähler-cone description
makes that choice a Kähler class.  After replacing $h$ by $-h$ if
necessary, it is the class of an ample bundle.  This sign change
does not change $V$ or $e$.  Fix the resulting polarized manifold
and let $D_h$ denote the component of the polarized period domain
containing its period.  Locally, the periods of its polarized
deformations fill an open subset of $D_h$.

\smallskip\noindent\emph{The algebraic family and the free arithmetic quotient.}

We need an algebraic family with an open marked period image.
Choose $k\in\Z_{>0}$ so that the bundle with class $kh$ on the preceding
manifold is very ample and has vanishing higher cohomology.  The
bundle extends locally over polarized deformations; bases of its
sections give nearby embeddings with the same Hilbert polynomial.
Locally, deformations of a complete-linear-system embedding are
deformations of the polarized pair together with projective frames
of these section spaces.  The pair deformation space is smooth by
local Torelli, and cohomology vanishing lets the frames vary freely.
There are no infinitesimal automorphisms, and finite automorphisms
act freely on frames.  A reduced irreducible Hilbert-scheme component
therefore has a smooth point with submersive polarized period map.

Restrict its universal family to a smooth connected Zariski-open
base $S$ of smooth irreducible holomorphic symplectic fibers,
retaining a submersive point.  This restriction is algebraic:
Hodge numbers are locally constant in a smooth projective family,
and the degeneracy locus of the relative two-form has closed image
in the base by properness.  Simple connectivity and differentiable
type are preserved in the connected smooth family.  The universal
hyperplane bundle has class $kh$.

Let
\[
 G_h(\Z)=\{g\in O(\Lambda,q):gh=h,\ gD_h=D_h\}.
\]
Choose a torsion-free finite-index congruence subgroup
$\Gamma\subset G_h(\Z)$ whose elements have determinant one.
Both groups act on the full lattice
$\Lambda$.  Restriction embeds $G_h(\Z)$ in $O(V,q)$,
because $\Lambda_{\Q}=\Q h\oplus V$; its image is arithmetic.
For example, a congruence subgroup acting trivially on the finite
gluing quotient between $\Z h\oplus(\Lambda\cap V)$ and $\Lambda$
extends across this quotient, which shows the required finite-index
property in the integral orthogonal group of $V$.

The family monodromy is contained in $G_h(\Z)$: it preserves the
BBF form and the class $kh$, hence also $h$, and preserves the
component of its periods.  Pass to the finite cover of $S$
corresponding to the inverse image of $\Gamma$ under this
monodromy representation.  A finite topological cover of a smooth
complex quasi-projective variety is algebraic and finite étale
\cite[Expos\'e~XII, Theorem~5.1]{SGA1}.
We keep one connected component and continue to call the base $S$.
This step requires only that $\Gamma$ has finite index in
$G_h(\Z)$; no hypothesis on the index of the family monodromy
image is used.

The action of $\Gamma$ on $D_h$ is free.  Indeed, the stabilizer
of a positive oriented two-plane is a compact subgroup of
$O(V_{\R},q)$.  Its intersection with a discrete integral subgroup
is finite and is therefore trivial in the torsion-free group
$\Gamma$.  This argument also removes any kernel of the projective
action.  Consequently
\[
 Q_h:=\Gamma\backslash D_h
\]
is a smooth arithmetic quotient and $D_h\to Q_h$ is a covering
map.  The Baily--Borel theorem \cite{BailyBorel} makes $Q_h$
quasi-projective.  Borel's algebraicity theorem
\cite[Theorem~3.10]{Borel} then makes the descended period map
\[
 \Phi:S\longrightarrow Q_h
\]
algebraic.  We have given the Hilbert-family construction directly;
compare \cite[Lemma~4.5]{Soldatenkov}.  At a chosen point $s_0\in S$, fix
the marking obtained by transport from the starting polarized
manifold and a local lift of $\Phi$ to $D_h$.  After shrinking an
analytic neighborhood, this lift is submersive and its image
contains an open set $U\subset D_h$ with the actual marking.
Freeness will later identify monodromy on the full second-cohomology
lattice from equality of these lifted periods.
\end{proof}

We now apply Lemma~\ref{lem:polarized-family-construction} to the
primitive class fixed at the start of this section. Its open marked
period image is the input for the next step.

\subsection{A rational ternary domain containing the fixed vector}

We now find a rational period subspace which meets $U$ and has
$\Q e$ as a cusp.  View $D_h$ as the space of positive oriented
real two-planes in
$V_{\R}$.  Let $W\in U$.  Write $e=e_W+e_-$ in the orthogonal
decomposition $V_{\R}=W\oplus W^\perp$.  Since $W^\perp$ is
negative definite, $e_W\ne0$; otherwise the nonzero vector $e$
would have negative square.  Moreover
\[
 q(e_-)=-q(e_W)<0.
\]
It follows that $W+\R e=W\oplus\R e_-$ has signature $(2,1)$.

For an explicit rational approximation, choose an oriented basis
$x,y$ of $W$ and rational vectors $x_j,y_j\in V$ tending to $x,y$.
For all sufficiently large $j$ the plane
$W_j=\langle x_j,y_j\rangle_{\R}$ is positive, has the chosen
orientation, and lies in $U$.  Nondegeneracy and signature are open
conditions on the Gram matrix, so
\[
 P_j=\langle e,x_j,y_j\rangle_{\Q}
\]
has signature $(2,1)$ for all sufficiently large $j$.  Fix such a
$P=P_j$.  The vector $e$ has remained fixed throughout this
approximation.  Let $D_P$ be the component of its
positive-plane domain that contains $W_j$.  Then $D_P\subset D_h$
and $D_P\cap U\ne\varnothing$.

An arithmetic subgroup of $SO(P,q)$ preserving $D_P$ can be made
to act on $\Lambda$ by extending it identically on $P^\perp$.
To see the integral issue explicitly, set
\[
 \Lambda_0=(P\cap\Lambda)\oplus(P^\perp\cap\Lambda).
\]
This is a finite-index sublattice of $\Lambda$.  Choose $N$ with
$N\Lambda\subset\Lambda_0$.  If an isometry $g$ of $P\cap\Lambda$
is congruent to the identity modulo $mN$, with $m\in\Z_{>0}$,
its extension satisfies
$(g-1)\Lambda\subset m\Lambda_0\subset m\Lambda$.
The same holds for its inverse, so the extension preserves
$\Lambda$; choosing $m$ sufficiently divisible places it in
$\Gamma$.  Take such a
finite-index subgroup $\Gamma_P$.  It is torsion-free and its
inclusion into $\Gamma$ is an inclusion of actual lattice
isometries.

The quotient
\[
 C_P=\Gamma_P\backslash D_P
\]
is a smooth arithmetic curve.  Its holomorphic map
$j:C_P\to Q_h$ is algebraic by Borel's theorem
\cite[Theorem~3.10]{Borel}.  Since $P$ contains
the rational null line $\Q e$, its smooth compactification
$\overline C_P$ has a cusp $c_e$ represented by that line.  We
always retain this chosen lift of the cusp in the boundary of
$D_P$, even if other translates represent the same point of
$\overline C_P$.

The construction also covers the two smallest possible ranks.
If $r=4$, then $V$ itself is a rational three-space of signature
$(2,1)$ and we simply take $P=V$.  Its noncompact arithmetic curve
has a cusp because the given $e$ is rational and isotropic.  If
$r=5$, the domain of signature $(2,2)$ is still a Hermitian
symmetric domain and the same ternary subdomain works.  General
constructions that obtain a null vector from Meyer's theorem need
an additional rank hypothesis; here $e$ is already supplied.

\subsection{An algebraic curve reaching the specified cusp}

It remains to reach this cusp in the actual polarized family and
identify its monodromy.  Form the algebraic fiber product
\[
 Z=S\times_{Q_h}C_P.
\]
Choose a point $z_0$ at which the marked period lift from $S$
equals the chosen point of $U\cap D_P$.  The
projection $Z\to C_P$ is submersive near $z_0$, because the local
period map from $S$ is submersive.  Thus the component through
$z_0$ dominates $C_P$.  It contains an irreducible algebraic curve
through a point in this selected local branch which dominates
$C_P$: intersect a quasi-projective embedding with general
hyperplanes through a smooth point, choosing their tangent
intersection transverse to the fiber of $Z\to C_P$, and take the
component through that point.  Normalize this curve and call it
$B$.  At the selected point the resulting map $g:B\to C_P$ is
nonconstant.  Pullback from $S$ gives a projective polarized family
over $B$ with the chosen marked local branch.

Take smooth projective completions.  The nonconstant map of curves
extends to a finite surjective morphism
\[
 \bar g:\overline B\longrightarrow\overline C_P.
\]
There is a finite set $F\subset\overline C_P$, containing all cusps,
branch values, and images of the finitely many points omitted
from $B$, such that over $C_P\setminus F$ the map is a finite
unramified covering and its source lies in $B$.  If necessary move
the chosen starting point a little within the same local branch
so that it maps outside $F$.

The following path choice fixes the marking at infinity.  Start
with its selected lift in $D_P$.  The inverse image in $D_P$ of
$F\cap C_P$ is a locally finite discrete set.  Its complement is
path connected.  Consequently we may take a path from the initial
lift towards the boundary point represented by the specific vector
$e$, avoiding that discrete set; sufficiently far along the path
it lies in a cusp neighborhood of $e$.  Project the path to
$C_P\setminus F$ and lift it through the covering from $B$, starting
at the chosen point.  The lift approaches some point
$b_e\in\overline B$ over $c_e$.  In a punctured disk around $b_e$
we therefore have the desired smooth algebraic polarized family.

Transport the original marking along this lifted path.  The
resulting family period lift to $D_h$ and the inclusion of the
chosen $D_P$-valued lift start at the same point and project to
the same map into $Q_h$.  Uniqueness of lifts for the covering
$D_h\to Q_h$ makes them equal all along the path and on the
universal cover of the resulting punctured disk.

Freeness now identifies the monodromy on all of $\Lambda$.
For a loop in $B$, let $M_B\in\Gamma$ be its
second-cohomology monodromy and let $\gamma_B\in\Gamma_P\subset
\Gamma$ be the deck transformation of its modular-curve lift,
using compatible conventions.  Equivariance and the preceding
equality of period lifts imply
\[
 M_B\widetilde\Phi(\widetilde b)
       =\gamma_B\widetilde\Phi(\widetilde b).
\]
Hence $\gamma_B^{-1}M_B$ fixes a point of $D_h$.  The action is
free, so $M_B=\gamma_B$ as isometries of the full lattice.
In particular, near the selected cusp the actual family
monodromy fixes our labelled $e$ and is the identity on $P^\perp$.
Thus the free action and the equality of the initial marked lifts
retain the precise class needed in the rotation argument.

\subsection{The exact period and monodromy formulas}

We have obtained a family whose marked periods lie in $D_P$ and
approach the cusp $\Q e$.  A rational hyperbolic basis makes both
their growth and their monodromy explicit.  Choose $f_0\in P$
with $q(e,f_0)=1$ and put
$f=f_0-\frac12q(f_0)e$.  Then $q(f)=0$ and $q(e,f)=1$.
The orthogonal complement of $\langle e,f\rangle$ in $P$ is a
positive rational line; choose a nonzero rational vector $v$ on
it and put $a=q(v)>0$.  This proves
\eqref{eq:family-hyperbolic-basis}.

For every positive period in $P_{\C}$ the pairing with $e$ is
nonzero.  Indeed, a positive real two-plane orthogonal to $e$
would put $e$ in its negative-definite orthogonal complement.
Normalize the period by $q(p,e)=1$.  Writing
$p=f+\tau v+\beta e$, its nullity gives
$\beta=-a\tau^2/2$.  Positivity becomes
\[
 q(p,\bar p)=a|\tau|^2-\frac a2(\tau^2+\bar\tau^2)
            =2a(\operatorname{Im}\tau)^2>0.
\]
Change the sign of $v$ so that the component $D_P$ is parametrized
by $\operatorname{Im}\tau>0$.  The point $\tau=\infty$ represents
the line $\Q e$.

For a rational number $u$, define an isometry $U_u$ of $P$ by
\[
 U_ue=e,\qquad U_uv=v-au e,\qquad
 U_uf=f+u v-\frac{au^2}{2}e.
\]
The Gram matrix in \eqref{eq:family-hyperbolic-basis} verifies
directly that $U_u$ preserves $q$, and substitution gives
\[
 U_up(\tau)=p(\tau+u).
\]
Taking $u$ to be a sufficiently divisible positive integer makes
this transformation belong to $\Gamma_P$.  The cusp stabilizer in the
torsion-free orientation-preserving arithmetic group is generated
by such a translation $\tau\mapsto\tau+w_0$, with $w_0>0$.
Here the width $w_0$ is rational because the basis $e,f,v$ is
rational and the generator preserves that rational structure.
The standard cusp coordinate is
\[
 u_c=\exp(2\pi i\tau/w_0).
\]
If a local parameter $t$ upstairs has ramification degree $b$ at
the cusp, then
\begin{equation}
 \label{eq:family-cusp-coordinate}
 u_c=t^b\beta(t),\qquad \beta(0)\ne0,
 \qquad
 \operatorname{Im}\tau
    =\frac{bw_0}{2\pi}(-\log|t|)+O(1).
\end{equation}
The error is uniform near the puncture.  A positive local loop
corresponds to $U_{bw_0}$.  The already established equality of
monodromy representations gives \eqref{eq:family-monodromy}
with $w=bw_0$.  Computing $(U_w-1)$ on the displayed basis gives
\eqref{eq:family-monodromy-image}.  Further finite base changes
replace $w$ by a positive integral multiple, so all the image
and nonvanishing statements survive.

Because only finitely many exceptional values were removed from
$\overline C_P$, a sufficiently small punctured cusp neighborhood
contains none of them.  In the preceding path construction we may
therefore choose its final lift in $D_P$ to be the ray
$\tau=\lambda+iy$ for any specified irrational $\lambda$ and all
large $y$.  The ray lifts to the chosen finite cover after its
initial branch is fixed.  It continues to lift after any further
finite base change.  The resulting path in the final parameter $t$
may be nonradial.

\subsection{Semistable reduction and the logarithmic form}

Take a projective algebraic model over a neighborhood of $b_e$ in
$\overline B$.  One-parameter semistable reduction in
characteristic zero, followed by restriction to a small disk,
gives, after finite base change and a projective modification
which is an isomorphism over the punctured base, a smooth total
space $\mathcal Y$ with reduced simple normal crossings central
fiber; see \cite[Chapter~II, Semi-stable Reduction Theorem,
pp.~53--54]{Semistable}.  The smooth fibers and their marked
periods remain those already constructed.  The polarization on
the punctured family extends to a line bundle on $\mathcal Y$:
choose a rational section, close its divisor in the smooth total
space, and use that every Weil divisor on a smooth variety is
Cartier.  Denote the extension by $\mathcal A$.  Its restriction
to each smooth fiber is ample; relative ampleness on the
semistable model is unnecessary.

The invariant vector $e$ is now a single-valued flat section of
$R^2\pi_*\Z$ on $\Delta^*$.  The pairing $q(-,e)$ is a
single-valued holomorphic functional on its associated flat
vector bundle.  Its restriction to the Hodge line is nowhere
zero, by \eqref{eq:family-period-identities}.  Therefore it
trivializes that Hodge line.  Holomorphic base change identifies
the Hodge line with the bundle of relative holomorphic two-forms,
so there is a unique holomorphic relative form $\sigma$ with
$q([\sigma],e)=1$.  On the universal cover its class is exactly
$p(\tau)$, including its scalar.  Thus $\eta=\sigma^n$ is a
single-valued holomorphic relative volume form on $\Delta^*$.

Polarizing the Fujiki identity \cite{Fujiki}, or simply comparing
the coefficient of $z^nw^n$ in its application to
$z[\sigma]+w[\bar\sigma]$, gives
\[
 \binom{2n}{n}\int_{Y_t}\sigma_t^n\wedge\bar\sigma_t^n
       =c_X2^n q(p(\tau),\overline{p(\tau)})^n.
\]
For the positive volume convention in $|\eta_t|^2$, any fixed
normalization factor can be absorbed in a constant.  Equations
\eqref{eq:family-period-identities} and
\eqref{eq:family-cusp-coordinate} consequently give
\[
 \int_{Y_t}|\eta_t|^2
       =c' (\operatorname{Im}\tau)^{2n}\asymp s^{2n}=s^d,
 \qquad c'>0,
\]
proving \eqref{eq:family-volume-growth} uniformly near the puncture.

Let $\mathcal K_{\log}$ denote the invertible sheaf
\[
 \det\!\left(
 \Omega^1_{\mathcal Y}(\log Y_0)
 \big/\pi^*\Omega^1_{\Delta}(\log\{0\})\right).
\]
The quotient is locally free of rank $d$ in the semistable
coordinates.  On $\Delta^*$ its determinant is the ordinary
relative canonical line, and $\eta$ is a section of it there.
We prove that $\eta$ extends holomorphically as a section of
$\mathcal K_{\log}$.

First consider a point of $Y_0$ lying on only one component.
There the family has product coordinates $(z_1,\ldots,z_d,t)$.
Write
\[
 \eta_t=F(z,t)\,\dd z_1\wedge\cdots\wedge\dd z_d.
\]
On a smaller product chart the submean inequality in the fiber
variables, followed by the global upper volume bound, gives
\[
 |F(z,t)|^2\le C\int_{Y_t}|\eta_t|^2\le C' s^d.
\]
For $m\ge1$ the negative Laurent coefficient in the $t$ variable
satisfies, uniformly on a smaller $z$-polydisk,
\[
 \left|\frac1{2\pi i}\int_{|t|=\rho}
                  F(z,t)t^{m-1}\,\dd t\right|
       \le C\rho^m|\log\rho|^{d/2}\longrightarrow0.
\]
All negative Laurent coefficients vanish.  Hence $F$ extends
holomorphically through $t=0$.  This extends $\eta$ over the
smooth locus of the reduced central divisor.  The crossing locus
has complex codimension at least two in the smooth total space.
Local trivializations of $\mathcal K_{\log}$ and Hartogs'
extension theorem then extend $\eta$ uniquely over that locus.
Thus the polynomial bound in $s$ suffices for logarithmic extension.

\subsection{Maximal growth forces a nonzero maximal residue}

We have the semistable family and the extended logarithmic form.
To apply Theorem~\ref{thm:torus}, we must still find a maximal
crossing with nonzero coefficient.  The lower growth bound forces
one, as follows.  At a point where precisely $k$ components of
$Y_0$ meet, take
coordinates with
\[
 t=z_0\cdots z_{k-1},\qquad 1\le k\le d+1.
\]
A local generator of $\mathcal K_{\log}$ restricts to the fiber as
\[
 \frac{\dd z_1}{z_1}\wedge\cdots\wedge
 \frac{\dd z_{k-1}}{z_{k-1}}
 \wedge\dd z_k\wedge\cdots\wedge\dd z_d.
\]
Its extended holomorphic coefficient is bounded on a smaller
chart.  In the variables $x_i=-\log|z_i|$, $1\le i\le k-1$,
the fiber part of a fixed coordinate polydisk has an expanding
simplex as its radial domain: each $x_i$ is bounded below, and
$\sum_i x_i\le s+O(1)$.  Its dimension is $k-1$.  Integrating
the angular variables and the remaining ordinary coordinates
therefore yields
\begin{equation}
 \label{eq:family-local-log-bound}
 \int_{Y_t\cap\text{chart}}|\eta_t|^2
          \le C(1+s^{k-1}).
\end{equation}

For $k=d+1$, the crossing is an isolated point $o$.  Suppose its
coefficient $F$ vanishes at $o$.  To improve
\eqref{eq:family-local-log-bound}, choose a polydisk
$|z_i|<\delta<1$, put $A=-\log\delta$, and write
\[
 x_i=A+(s-(d+1)A)u_i,\qquad
 u_i>0,\quad \sum_{i=1}^d u_i<1.
\]
For almost every $u$ in this fixed simplex, all $d+1$ coordinates
$z_0,\ldots,z_d$ tend to zero as $s\to\infty$, uniformly in
their arguments.  The rescaled integrand $|F|^2$ thus tends to
$|F(o)|^2=0$ almost everywhere.  It is uniformly bounded and the
simplex and angle torus have finite measure.  Dominated convergence
gives
\begin{equation}
 \label{eq:family-zero-residue}
 \int_{Y_t\cap\text{chart}}|\eta_t|^2=o(s^d)
 \quad\text{if }F(o)=0.
\end{equation}

Properness gives finitely many such charts covering the central
fiber and, after shrinking the base, all nearby fibers.  If there
were no maximal crossings, every chart would contribute
$O(s^{d-1})$.  If maximal crossings existed but all their
coefficients vanished, their contributions would be $o(s^d)$
by \eqref{eq:family-zero-residue}, with all other contributions
again $O(s^{d-1})$.  Either conclusion contradicts the lower bound
in \eqref{eq:family-volume-growth}.  At least one maximal crossing
therefore has nonzero coefficient.  This proves
\eqref{eq:family-maximal-residue}.  At that point the functions
$z_1,\ldots,z_d$ may be chosen to be algebraic local defining
functions of the corresponding components; the unit in the
product equation is absorbed in $z_0$.  They are thus restrictions
of rational functions on the algebraic total space, as required
for the intersection estimates in Theorem~\ref{thm:torus}.
This completes the proof of
Proposition~\ref{prop:family}.

\begin{corollary}
\label{cor:family-torus}
Fix an irrational $\lambda$ and choose the lifted path in
Proposition~\ref{prop:family} with $\tau=\lambda+iy$.
For every sufficiently large $y$, the corresponding smooth fiber
contains an embedded special Lagrangian real $d$-torus for its
Ricci-flat metric in the polarization class.  This torus is
isotopic to a coordinate angle torus in the chart
\eqref{eq:family-maximal-residue}.  The period and monodromy on
that fiber are the exact marked data
\eqref{eq:family-period}--\eqref{eq:family-monodromy-image}.
\end{corollary}

\begin{proof}
The algebraicity, projectivity, extension of the polarization,
volume growth, logarithmic extension, and nonzero coefficient
required by Theorem~\ref{thm:torus} are all supplied by
Proposition~\ref{prop:family}.  The theorem applies to every
sufficiently small parameter, in particular along the chosen path,
where $\operatorname{Re}\tau=\lambda$ and the marking is already fixed.
\end{proof}
